\section{Puentes}
\label{sec:puentes}

Las cadenas de bloques descritas en las secciones anteriores mantienen cada una su propio estado
y su propio mecanismo de consenso, y ninguna de ellas puede, por sí sola, conocer lo que ocurre en
otra. Un contrato desplegado en Optimism no tiene acceso al estado de Ethereum más allá de los
mecanismos nativos que el propio OP Stack define (Sección~\ref{sec:op_puente_nativo}), y entre
cadenas independientes no existe vínculo alguno. Los \textbf{puentes} (\textit{bridges}) son los
sistemas que permiten trasladar activos e información entre cadenas en ausencia de ese vínculo.
Esta sección presenta los modelos con que se construyen, los incidentes de seguridad que han
sufrido y las propuestas que sustituyen la confianza en intermediarios por pruebas de
conocimiento cero.

\subsection{Interoperabilidad y modelos de puente}
\label{sec:puentes_modelos}

Se entiende por \textbf{interoperabilidad} entre cadenas de bloques la capacidad de un sistema de
transferir activos o datos de una cadena a otra de manera que la cadena receptora pueda actuar
sobre ellos con garantías comparables a las que ofrece la cadena de origen~\cite{belchior2021survey}.
Un puente materializa esta capacidad mediante tres componentes genéricos. En la cadena de origen,
un contrato recibe la operación del usuario y la registra. En la cadena de destino, otro contrato
ejecuta la operación correspondiente una vez que se convence de que la primera ocurrió. Entre
ambas, uno o varios agentes fuera de la cadena observan la cadena de origen y transmiten a la de
destino la información necesaria, puesto que ninguna cadena puede iniciar por sí misma una
comunicación con otra~\cite{zhang2023bridgesok}.

\paragraph{Modelos de transferencia de activos.} Un activo nativo de una cadena no puede
trasladarse literalmente a otra, ya que su existencia se reduce a un saldo registrado en el estado
de la primera. Lo que un puente transfiere es, en rigor, un derecho, y la forma de representarlo
en la cadena de destino da lugar a tres modelos, resumidos en la
Tabla~\ref{tab:puentes_transferencia}. En el modelo de \textit{bloqueo y acuñación}
(\textit{lock-and-mint}), el usuario deposita los tokens en un contrato de custodia de la cadena
de origen y el contrato de destino acuña una cantidad equivalente de un token representativo,
denominado habitualmente \textit{token envuelto} (\textit{wrapped token}). El retorno sigue el
camino inverso, de modo que el token representativo se quema en destino y los tokens originales
se liberan de la custodia. En el modelo de \textit{quema y acuñación} (\textit{burn-and-mint}),
aplicable cuando un mismo emisor controla el token en ambas cadenas, los tokens se destruyen en
origen y se emiten en destino, sin que exista una reserva bloqueada. El protocolo CCTP de Circle,
que traslada la moneda estable USDC entre cadenas a partir de una atestación emitida por el propio
emisor, constituye un ejemplo de este esquema~\cite{circlecctp}. Por último, en el modelo de
\textit{bloqueo y liberación} (\textit{lock-and-unlock}), el puente mantiene reservas de liquidez
del activo nativo en ambas cadenas, y la transferencia consiste en recibir el depósito en una de
ellas y entregar fondos de la reserva en la otra.

\begin{table}[htbp]
	\centering
	\caption{Modelos de transferencia de activos entre cadenas.}
	\label{tab:puentes_transferencia}
	\small
	\begin{tabular}{>{\raggedright\arraybackslash}p{2.7cm} >{\raggedright\arraybackslash}p{4.2cm} >{\raggedright\arraybackslash}p{3.3cm} >{\raggedright\arraybackslash}p{3.6cm}}
		\toprule
		\textbf{Modelo} & \textbf{Operación en origen y en destino} & \textbf{Activo recibido} &
		\textbf{Requisito principal} \\
		\midrule
		Bloqueo y acuñación & Custodia en origen y acuñación en destino & Token representativo &
		La custodia debe respaldar en todo momento lo acuñado \\
		\addlinespace
		Quema y acuñación & Destrucción en origen y emisión en destino & Token nativo del mismo
		emisor & El emisor debe controlar el token en ambas cadenas \\
		\addlinespace
		Bloqueo y liberación & Depósito en origen y entrega desde una reserva en destino & Token
		nativo de la cadena de destino & Reservas de liquidez suficientes en ambas cadenas \\
		\bottomrule
	\end{tabular}
\end{table}

Los tres modelos comparten una misma pregunta, que constituye el núcleo del problema de seguridad
de cualquier puente. El contrato de destino debe decidir si la operación declarada en la cadena
de origen ocurrió realmente, y una decisión errónea en este punto permite crear tokens sin
respaldo o retirar fondos de una reserva sin haber depositado nada. En el modelo de bloqueo y
acuñación, que es el adoptado en este trabajo, cada token acuñado sin el bloqueo correspondiente
deja la custodia insuficiente para respaldar a todos los tenedores del token representativo. La
Figura~\ref{fig:puente_confianza} ilustra este modelo en su forma más difundida, en la que un
\textit{relayer} de confianza informa los bloqueos y el contrato de destino acuña sin verificar
nada más que la identidad de quien lo invoca.

\begin{figure}[htbp]
	\centering
	\begin{tikzpicture}[
		node distance=0.6cm and 1.4cm,
		caja/.style={draw, rounded corners=2pt, align=center, minimum width=3.6cm,
			minimum height=1cm, font=\small},
		cadena/.style={draw, dashed, rounded corners=4pt, inner sep=0.3cm},
		flecha/.style={-{Stealth[length=2.2mm]}, thick},
		dato/.style={font=\footnotesize, align=center}
		]
		\node[caja] (usuario) {Usuario};
		\node[caja, below=1.4cm of usuario] (custodia) {Contrato de custodia\\(tokens bloqueados)};
		\node[caja, right=2.6cm of usuario, fill=gray!15] (relayer) {\textit{Relayer} de confianza\\(clave privada)};
		\node[caja, right=2.6cm of custodia] (acunador) {Contrato de destino\\(acuña el token
			representativo)};

		\node[cadena, fit=(custodia), label={[font=\footnotesize]below:Cadena de origen}] {};
		\node[cadena, fit=(acunador), label={[font=\footnotesize]below:Cadena de destino}] {};

		\draw[flecha] (usuario) -- node[left, dato] {1. bloqueo} (custodia);
		\draw[flecha] (custodia) -- node[above, dato, sloped] {2. evento} (relayer);
		\draw[flecha] (relayer) -- node[right, dato] {3. orden firmada\\de acuñación} (acunador);
	\end{tikzpicture}
	\caption[Puente de bloqueo y acuñación con un \textit{relayer} de confianza.]{Puente de bloqueo
	y acuñación con un \textit{relayer} de confianza. El contrato de destino acepta cualquier orden
	firmada con la clave del \textit{relayer}, por lo que la seguridad de los fondos depende
	íntegramente de esa clave y de la honestidad de quien la controla.}
	\label{fig:puente_confianza}
\end{figure}

\paragraph{Modelos de verificación.} La forma en que el contrato de destino responde a la
pregunta anterior determina el supuesto de confianza del puente y permite clasificar los diseños
existentes en cuatro familias~\cite{zhang2023bridgesok, belchior2021survey}, que se comparan en la
Tabla~\ref{tab:puentes_verificacion}.

En los puentes de \textbf{verificación externa}, un conjunto de agentes ajenos a ambas cadenas
atestigua los eventos de origen y el contrato de destino se limita a comprobar sus firmas. El
conjunto puede reducirse a un único \textit{relayer}, como en la Figura~\ref{fig:puente_confianza},
o ampliarse a una multifirma que exige un umbral de firmas entre varios participantes, o a un comité que firma mediante
computación multiparte. Estos diseños son económicos y rápidos, pero su seguridad equivale a la
de las claves del conjunto, y un adversario que obtenga el umbral de firmas puede acuñar o
liberar fondos sin que la cadena de destino pueda advertirlo.

Los puentes de \textbf{verificación optimista} aceptan provisionalmente las afirmaciones de un
agente y abren un período durante el cual cualquier vigilante (\textit{watcher}) puede impugnarlas
con evidencia de fraude. Su supuesto de confianza se reduce a la existencia de al menos un
vigilante honesto y activo durante cada período de impugnación, a cambio de una latencia que suele
medirse en minutos u horas. Se trata del mismo principio que emplean los \textit{rollups}
optimistas descritos en la Sección~\ref{sec:op_fraude}, trasladado a la comunicación entre cadenas
independientes.

Los puentes de \textbf{verificación nativa} ejecutan en la cadena de destino un cliente ligero de
la cadena de origen, que comprueba las firmas de consenso de cada encabezado y, a partir de él, las
pruebas de inclusión de cada operación. El protocolo IBC del ecosistema Cosmos es el ejemplo más
extendido de este enfoque, y su especificación define cómo cada cadena mantiene clientes ligeros
de sus contrapartes~\cite{ibcics02}. Con este diseño, el puente no añade supuestos de confianza
a los de las propias cadenas, pero solo resulta viable cuando la verificación del consenso de
origen es asequible en destino. Para Ethereum no lo es, ya que exigiría verificar firmas BLS
agregadas sobre la curva BLS12-381 y cientos de hashes en cada actualización
(Sección~\ref{sec:eth_cadena_verificacion}).

Los puentes de \textbf{verificación con pruebas de conocimiento cero} resuelven esta limitación
trasladando la ejecución del cliente ligero fuera de la cadena. Un probador ejecuta la
verificación del consenso y de la inclusión y genera una prueba sucinta de que la ejecutó
correctamente, mientras que el contrato de destino verifica únicamente esa prueba con un costo
acotado~\cite{xie2022zkbridge}. El supuesto de confianza coincide con el del cliente ligero
verificado, al que se suma la solidez del sistema de pruebas, y el agente que transmite la
información deja de tener autoridad sobre el resultado, puesto que una prueba de un enunciado
falso es rechazada sin importar quién la presente.

\begin{table}[htbp]
	\centering
	\caption{Modelos de verificación de puentes entre cadenas.}
	\label{tab:puentes_verificacion}
	\small
	\begin{tabular}{>{\raggedright\arraybackslash}p{2.6cm} >{\raggedright\arraybackslash}p{4.4cm} >{\raggedright\arraybackslash}p{2.6cm} >{\raggedright\arraybackslash}p{4.2cm}}
		\toprule
		\textbf{Modelo} & \textbf{Supuesto de confianza} & \textbf{Latencia} &
		\textbf{Costo en destino} \\
		\midrule
		Externa & Honestidad de al menos el umbral de firmantes & Baja & Bajo, comprobación de firmas
		convencionales \\
		\addlinespace
		Optimista & Al menos un vigilante honesto y activo & Alta, por el período de impugnación &
		Bajo en ausencia de disputas \\
		\addlinespace
		Nativa & Seguridad del consenso de origen & Baja & Alto, ejecución completa del cliente
		ligero \\
		\addlinespace
		Conocimiento cero & Seguridad del consenso de origen y solidez del sistema de pruebas &
		Media, por la generación de la prueba & Moderado, verificación de una prueba sucinta \\
		\bottomrule
	\end{tabular}
\end{table}

\subsection{Incidentes de seguridad}
\label{sec:puentes_incidentes}

Los puentes concentran en sus contratos de custodia el respaldo de todos los tokens
representativos que han emitido, lo que los convierte en objetivos de gran valor. Según
Chainalysis, los ataques a puentes representaron el 69\,\% de los fondos sustraídos en el
ecosistema de criptoactivos durante 2022 hasta el mes de agosto, con pérdidas cercanas a dos mil
millones de dólares estadounidenses distribuidas en trece incidentes~\cite{chainalysis2022bridges}.
La Tabla~\ref{tab:puentes_incidentes} reúne cinco de los casos de mayor magnitud de ese año, con
la causa documentada por los propios equipos afectados o por los analistas que los investigaron.

\begin{table}[htbp]
	\centering
	\caption{Incidentes de seguridad relevantes en puentes entre cadenas durante 2022.}
	\label{tab:puentes_incidentes}
	\small
	\begin{tabular}{>{\raggedright\arraybackslash}p{2.6cm} >{\raggedright\arraybackslash}p{1.5cm} >{\raggedright\arraybackslash}p{2.8cm} >{\raggedright\arraybackslash}p{4.6cm} >{\raggedright\arraybackslash}p{1.8cm}}
		\toprule
		\textbf{Puente} & \textbf{Fecha} & \textbf{Pérdida aproximada} & \textbf{Causa} &
		\textbf{Modelo} \\
		\midrule
		Ronin~\cite{roninpostmortem} & Marzo & 173\,600 ETH y 25,5 millones de USDC (unos US\$\,620
		millones) & Compromiso de cinco de las nueve claves de los validadores del puente & Externa \\
		\addlinespace
		Wormhole~\cite{wormholeincident} & Febrero & 120\,000 wETH (unos US\$\,325 millones) &
		Omisión de la verificación de firmas de los guardianes en el contrato de Solana & Externa \\
		\addlinespace
		Harmony Horizon~\cite{harmonyhorizon} & Junio & Unos US\$\,100 millones & Robo de dos claves
		de una multifirma de dos de cinco participantes & Externa \\
		\addlinespace
		Nomad~\cite{nomadrca} & Agosto & Unos US\$\,186 millones & Una actualización dejó como raíz
		de confianza el valor cero, con lo que se aceptaron mensajes no probados & Optimista \\
		\addlinespace
		BNB Token Hub~\cite{verichainsbnb} & Octubre & 2 millones de BNB acuñados (unos US\$\,570
		millones) & Falsificación de una prueba de Merkle IAVL aceptada por el verificador & Nativa \\
		\bottomrule
	\end{tabular}
\end{table}

Los casos reunidos admiten dos lecturas complementarias. Los incidentes de Ronin y Harmony
Horizon no explotaron ningún error de programación, sino que el adversario obtuvo las claves
suficientes para firmar órdenes que el contrato de destino consideraba legítimas por
definición~\cite{roninpostmortem, harmonyhorizon}. Se trata de la falla característica de la
verificación externa, en la que la seguridad de los fondos se reduce a la custodia de unas pocas
claves privadas, y es la que motiva la sustitución de la confianza en firmantes por verificación
criptográfica del estado de origen.

Los incidentes de Wormhole, Nomad y BNB Token Hub responden a una causa distinta. En los tres, la
lógica de verificación existía, pero su implementación aceptaba evidencia que no debía aceptar.
En Wormhole, el contrato no comprobó la autenticidad de la cuenta de la que leía las firmas
verificadas, de modo que un mensaje falsificado pasó por válido~\cite{wormholeincident}. En
Nomad, una actualización inicializó la raíz de confianza con el valor cero, que coincide con el
valor por omisión de los mensajes nunca probados, y cualquier mensaje pasó a considerarse
probado~\cite{nomadrca}. En BNB Token Hub, un defecto en la verificación de las pruebas de
Merkle del árbol IAVL permitió construir una prueba falsa de un depósito
inexistente~\cite{verichainsbnb}. Este último caso resulta especialmente ilustrativo, pues el
puente verificaba pruebas de inclusión respecto del estado de la cadena de origen y aun así fue
vulnerado. De ello se desprende que reemplazar a los firmantes por pruebas criptográficas elimina
la primera clase de fallas, pero no la segunda. La seguridad del puente depende también de que el
enunciado probado capture exactamente la condición de acuñación, de que el verificador esté
correctamente implementado y de que el contrato receptor impida la reutilización de una misma
prueba, consideraciones que ya se anticiparon en el Capítulo~\ref{chp:conocimiento_cero} y que se
retoman al analizar el diseño propuesto en la Sección~\ref{sec:puente_aceptacion}.

\subsection{Estado del arte}
\label{sec:puentes_estado_arte}

La propuesta que estableció la viabilidad práctica de los puentes verificados con pruebas de
conocimiento cero es zkBridge, de Xie et al.~\cite{xie2022zkbridge}. Su arquitectura separa un
contrato actualizador, que incorpora encabezados de la cadena de origen solo cuando se acompañan
de una prueba de que satisfacen la regla de validación del cliente ligero, de los contratos de
aplicación, que verifican sobre esos encabezados las pruebas de inclusión de cada operación. Para
que la generación de las pruebas resulte viable, los autores introdujeron deVirgo, una versión
distribuida del sistema de pruebas Virgo que reparte el circuito entre varias máquinas, y una
segunda capa que demuestra con Groth16 la validez de la prueba anterior, a fin de reducir el
costo de verificación en destino. El análisis detallado de esta propuesta y su relación con el
diseño de este trabajo se presentan en la Sección~\ref{sec:zkbridge_berkeley}.

Polyhedra Network llevó esta propuesta a un despliegue comercial, con clientes ligeros de varias
cadenas verificados mediante la misma combinación de deVirgo y una capa recursiva basada en
Groth16. Según sus autores, la capa recursiva reduce el costo de verificación en Ethereum desde
unos 80 millones de unidades de gas hasta unas 227\,000~\cite{polyhedrazkbridge}. Al tratarse de
cifras informadas por el propio equipo de desarrollo, se consignan como referencia y no como
resultado verificado de forma independiente.

Succinct Labs desarrolló Telepathy, un cliente ligero de Ethereum para otras cadenas que verifica
dentro de un circuito las firmas del comité de sincronización descrito en la
Sección~\ref{sec:eth_sync_committee}, en lugar de la finalidad del conjunto completo de
validadores~\cite{succinctproofconsensus, telepathydocs}. Telepathy separa dos operaciones. La
primera prueba que una supermayoría del comité vigente firmó un nuevo encabezado, y la segunda
prueba que el comité del período siguiente fue seleccionado correctamente, lo que permite al
contrato rotar el comité sin depender de su propietario~\cite{telepathydocs}. El sistema se empleó
para dar seguridad al puente entre Ethereum y Gnosis Chain~\cite{succinctproofconsensus}.

Con la aparición de las máquinas virtuales de conocimiento cero, el mismo equipo publicó SP1
Helios, que reemplaza el circuito escrito a mano por un programa en Rust ejecutado en la zkVM SP1
(Sección~\ref{sec:sp1}). El programa reutiliza la biblioteca de cliente ligero Helios para
comprobar que el encabezado de consenso fue firmado por el comité de sincronización y que el
encabezado de ejecución contenido en él es correcto, y el contrato de destino verifica la prueba
resultante~\cite{sp1helios}. Este cambio ilustra la alternativa que se estudia en el presente
trabajo, en la que el enunciado deja de expresarse como un circuito en un lenguaje específico y
pasa a expresarse como un programa convencional cuya ejecución se demuestra.

Las propuestas reseñadas comparten el supuesto de confianza en el comité de sincronización, que
es también el adoptado en este trabajo, y se concentran en la verificación del consenso, dejando
la prueba de inclusión de cada operación a cargo de las aplicaciones. El presente trabajo se
diferencia de ellas en tres aspectos. En primer lugar, demuestra en un mismo puente el consenso y
la inclusión de la operación de bloqueo, incluida la ejecución exitosa registrada en su recibo.
En segundo lugar, implementa los mismos enunciados con dos herramientas de naturaleza distinta,
el lenguaje de circuitos Noir y la zkVM SP1, lo que permite compararlas sobre un mismo problema.
En tercer lugar, orienta la verificación hacia una red de segunda capa, Optimism Sepolia, cuyas
particularidades de costo se describieron en la Sección~\ref{sec:op_tarifas}. El diseño
resultante se expone en el Capítulo~\ref{chp:zkbridge}.
